% 来源：OEBPS/Text/01_26.xhtml（OPF spine 顺序；仅第一册）
\cleardoublepage
\BookEntry[戴维·S.兰德斯《国富国穷》]{戴维·S.兰德斯《国富国穷》}

\BookCover{images/书单12国富国穷.jpg}{[美] 戴维·S.兰德斯著| 门洪华译| 新华出版社2010年版}

\begin{BookSummaryBox}
\textbf{这本书回答的正是：为什么一些国家富裕，而另一些国家贫穷？}
\par
\end{BookSummaryBox}

\subsubsection*{这本书谈什么？}

这本书回答了今天人们最感兴趣、同时也争议最大的一个问题：为什么一些国家成功实现了经济发展，而另一些国家却深陷贫穷的泥淖？
\par

作者批判、发展了斯密《国富论》的基本观点，将地理、历史、科技、哲学、宗教、国际政治、文化等因素综合在一起，多层次地分析了财富创造、国家兴衰更替的根本原因，提出了人类应该吸取历史教训，创造性地走出一条人、自然与社会协调发展的新路。
\par

这本书因其主旨宏大，思想精深，而被誉为“新《国富论》”。
\par

\subsubsection*{芒格为什么谈起这本书？}

在伯克希尔·哈撒韦1999财年的年会上，芒格说：
\par

\begin{quotation}
这本书重要的书（《国富国穷》）考察的是，随着时间的推移，财富与文化会如何相互作用，两者又会对国家产生怎样的影响。一些文化区虽然具有一定优势，却从来没有成为世界强权，这本书对这一现象也给出了引人深思的理性解释。这本书还检视了不同国家兴衰各异的原因。这些原因是每一个人都应该了解的。这样我们就可以从这些挑战中学习，而不再重复那些悲惨的命运。
\par
\end{quotation}

\subsubsection*{《国富国穷》作者戴维·S.兰德斯如是说}

“相对于善心而言，我更钟爱事实。”
\par

“荷兰人的成功，集中反映了他们在工作和经商过程中所具有的、在龟兔赛跑的寓言故事中乌龟所体现的精神。丰厚的战利品和奖赏当然好，但以长远目光来看（千万不要忽视长期效益），最重要的还是从那些不起眼的、风险较低的小宗生意中获取的利润，它们会积少成多，而且从来都不会令人失望。”
\par
